% !TeX root = ../初等数学笔记.tex
% 由主文档通过 \input 载入；请编译主文档。
\setcounter{chapter}{6} % 保留原书章号。
\chapter{函数与图像}
\label{chap:em-functions}
方程关注哪些输入产生指定输出，函数则研究全部输入与输出之间的关系。同一关系可以用式子、图像和语言表达；这些表示应彼此解释，而不能用不精确的图形观察替代定义。

\section{函数、定义域与值域}
\begin{definition}
\emph{函数}（function）$f:D\to Y$ 给每个 $x\in D$ 指定唯一元素 $f(x)\in Y$。$D$ 称为定义域，$Y$ 称为陪域，集合 $f(D)=\{f(x):x\in D\}$ 称为值域。
\end{definition}
函数不必有一个统一解析式。例如分段计费可由几段规则共同规定。相同解析式配上不同定义域也能产生不同函数：$x^2$ 在 $\mathbb R$ 上值域为 $[0,+\infty)$，在 $[1,2]$ 上值域为 $[1,4]$。定义域限制来自题目条件与运算合法性二者的交集。

若规定函数相等还须具有相同陪域，就应把陪域也视为函数数据；本书讨论实值函数时统一把陪域取为 $\mathbb R$，重点比较定义域和对应规则。求值域需要说明每个候选输出确实能由某个允许输入取得，而非只给出一个包含输出的区间。

\section{图像、坐标与基本函数}
函数图像是点集 $\{(x,f(x)):x\in D\}$。每条竖直直线至多与它相交一次，这是单值性的几何表达。一次函数 $kx+b$ 的图像是直线；$k=0$ 时为水平直线。二次函数配方为
\[
 ax^2+bx+c=a\left(x+\frac b{2a}\right)^2+c-\frac{b^2}{4a},\qquad a\ne0,
\]
由平方非负即可确定顶点、对称轴和最值，不必预先使用导数。

反比例函数 $k/x$ 在 $k\ne0$ 时排除 $x=0$，取值也不为零；其两支的位置由 $k$ 的符号确定。$|x|$ 由 $x$ 与 $-x$ 两段组成。作图时先确定定义域和关键点，再研究区间内性质，不能只把少数采样点用任意光滑曲线连接。

\section{单调性、奇偶性与有界性}
在区间 $I$ 上，若任取 $x_1<x_2$ 都有 $f(x_1)<f(x_2)$，称 $f$ 严格递增；用 $\le$ 替换 $<$ 得到非严格递增。递减相应定义。例如在 $[0,+\infty)$ 上，
\[
 x_2^2-x_1^2=(x_2-x_1)(x_2+x_1)>0,
\]
故平方函数严格递增；在负半轴上方向相反。$1/x$ 在正、负半轴上分别严格递减，却不在整个不连通定义域上严格递减，例如 $-1<1$ 而 $-1<1$ 的函数值顺序没有反转。

定义域关于原点对称时，$f(-x)=f(x)$ 称为偶函数，$f(-x)=-f(x)$ 称为奇函数。奇函数只有在 $0$ 属于定义域时才能推出 $f(0)=0$。若存在 $M$ 使所有 $f(x)\le M$，称 $M$ 为上界；只有某个输入实际取得这个上界且它不小于任何函数值时，才称它为最大值。$f(x)=x$ 在 $(0,1)$ 上有界，却没有最大值或最小值。

\section{复合函数、反函数与图像变换}
复合函数 $(f\circ g)(x)=f(g(x))$ 的定义域是 $\{x\in D_g:g(x)\in D_f\}$。例如 $f(u)=\sqrt u$、$g(x)=x^2-1$ 的复合定义域为 $(-\infty,-1]\cup[1,+\infty)$。次序通常不能交换。

若 $f$ 在 $D$ 上为单射，即相同输出只能来自相同输入，则对每个 $y\in f(D)$ 可以唯一恢复 $x$，从而定义\emph{反函数}（inverse function）$f^{-1}:f(D)\to D$。平方函数在整个实轴上不单射，限制在 $[0,+\infty)$ 后反函数为 $\sqrt x$；限制在 $(-\infty,0]$ 后则为 $-\sqrt x$。反函数的图像是原图像关于直线 $y=x$ 的反射，$f^{-1}$ 不是倒数 $1/f$。

图像变换可以从点对应推导：$(x,f(x))$ 在水平右移 $h$ 后成为 $(x+h,f(x))$，故新式为 $y=f(x-h)$；竖直平移得到 $f(x)+k$，纵向缩放得到 $Af(x)$。$|f(x)|$ 把横轴下方图像翻到上方，$f(|x|)$ 则把原右半图像向左镜像。

\section{零点、交点与函数观点下的求解}
$f(x)=g(x)$ 等价于公共定义域内 $f-g$ 的零点问题；$f(x)>g(x)$ 表示前者图像位于后者上方。严格单调函数至多有一个零点，因为两个不同输入不能得到同一输出，但单调性不保证零点存在，例如 $2^x>0$，其性质见下一章。

常用的零点存在依据是连续函数的介值性质：若 $f$ 在 $[a,b]$ 连续且 $f(a)f(b)<0$，则 $(a,b)$ 内有零点。连续性意为输入充分接近时输出也充分接近；其严格定义及介值定理证明属于分析基础，本书使用时明确列出条件。多项式连续，因此可用端点符号与单调性共同证明它在某区间恰有一根，图像只用于帮助定位。

\section{例题、习题与学习检查}
\begin{example}
求 $f(x)=x^2-4x+5$ 在 $[0,3]$ 上的值域。
\end{example}
\begin{solve}
$f(x)=(x-2)^2+1$。最小值 $1$ 在 $x=2$ 取得；区间中离 $2$ 最远的点为 $0$，最大值为 $5$。对每个 $y\in[1,5]$，取 $x=2-\sqrt{y-1}\in[0,2]$ 即有 $f(x)=y$，故值域恰为 $[1,5]$；这一构造补足了“所有中间值均能取到”的论证。
\end{solve}
\begin{enumerate}
\item 求 $f(x)=1/(x-2)$ 的定义域和值域。答：分别为 $\mathbb R\setminus\{2\}$ 与 $\mathbb R\setminus\{0\}$。
\item 求 $f(x)=3x-2$ 的反函数。答：$f^{-1}(x)=(x+2)/3$。
\item 证明 $x^3+x=1$ 恰有一个实根。提示：作差证明严格递增，再在 $[0,1]$ 上使用介值性质。
\end{enumerate}

